% Related Work
\section{Background and Related Work}
\label{sec:related_work}

\subsection{Multi-Agent Systems for Software Engineering}

Multi-agent approaches to software engineering have gained momentum with the
advancement of LLM capabilities. We survey representative systems and frameworks:

\paragraph{ChatDev} \citep{qian2023chatdev} simulates a virtual software company
with role-playing agents (CEO, CTO, Programmer, Tester) that communicate through
structured chat. While effective for generating simple applications, the sequential
communication pattern creates bottlenecks, since each agent must wait for prior agents
to complete their turns before contributing.

\paragraph{MetaGPT} \citep{hong2023metagpt} introduces standardized operating
procedures (SOPs) that structure agent collaboration around software artifacts
(PRDs, design documents, code). This reduces communication overhead compared to
free-form chat but still relies on sequential handoffs between roles.

\paragraph{Multi-Agent Debate (MAD)} \citep{liang2023encouraging} uses multiple
agents that debate and refine solutions iteratively. While debate can improve
solution quality through diverse perspectives, the synchronous turn-taking
limits parallelism.

\paragraph{Autonomous Agent Frameworks} AutoGPT \citep{autogpt2023} and
BabyAGI \citep{babyagi2023} pioneered recursive self-directed goal decomposition
but lack dependency management and formal quality gates, yielding unpredictable
execution paths. CrewAI \citep{crewai2024} coordinates role-defined personas
through sequential pipelines without built-in validation or parallel dispatch,
while LangGraph \citep{langgraph2024} models workflows as graphs at the level of
individual LLM invocations. SPOQ instead operates at task granularity (1--4 hour
atomic units) with enforced quality thresholds between waves.

\medskip
SPOQ differs from these systems in three ways: (1) wave-based parallel execution
rather than sequential role-playing; (2) explicit validation gates with quantified
metrics; and (3) structured human integration rather than full autonomy.

\subsection{Task Decomposition and Dependency Management}

Hierarchical task decomposition has roots in classical AI planning \citep{nau2003shop2}.
Modern approaches apply LLMs to generate task breakdowns:

\paragraph{Hierarchical Task Networks (HTNs).} Classical HTN planners decompose
abstract tasks into primitive actions with ordering constraints. SPOQ adapts
this structure for multi-agent software engineering.

\paragraph{DAG Scheduling.} Topological sorting of directed acyclic graphs is
well-established for parallel task scheduling \citep{coffman1972optimal}. SPOQ
applies these algorithms to LLM-generated task dependencies, computing wave
assignments that maximize parallelism while respecting precedence constraints.

\paragraph{Critical Path Analysis.} PERT and CPM methods identify the longest
path through a dependency graph \citep{kelley1959critical}. SPOQ uses critical
path analysis to estimate minimum wall-clock execution time and identify
bottleneck tasks.

\newpage
\subsection{Human-AI Collaboration in AI-Native Engineering}

Human-in-the-loop (HITL) approaches to AI systems span a spectrum from passive
oversight to active collaboration:

\paragraph{Prompt Engineering.} Users craft prompts to guide LLM behavior but
have limited ability to influence intermediate reasoning or correct course
mid-execution \citep{white2023prompt}.

\paragraph{AI Pair Programming.} Tools like GitHub Copilot \citep{chen2021codex}
suggest code completions while humans retain editing control. This inverts SPOQ's
model: humans do the work while AI assists, rather than agents doing the work
while humans validate.

\paragraph{Supervised Autonomy.} Some systems allow human intervention at
checkpoints \citep{wu2023autogen}. SPOQ extends this with bidirectional
communication: not only can humans intervene, but agents can explicitly
request human assistance when facing ambiguity.

\subsection{Autonomous Agent Frameworks}

Early autonomous agent frameworks pioneered recursive goal decomposition and
self-directed task execution, establishing foundational patterns for agentic AI:

\paragraph{AutoGPT and BabyAGI.}
AutoGPT \citep{autogpt2023} and BabyAGI \citep{babyagi2023} introduced recursive
goal decomposition where agents autonomously break objectives into subtasks and
execute them. These early systems demonstrated that LLMs could function as
autonomous agents but lack explicit dependency management and formal quality
gates, leading to unpredictable execution paths. SPOQ addresses these limitations
through DAG-based scheduling with dual validation gates.

\subsection{Autonomous Software Engineering Agents}

A new category of autonomous software engineering agents has emerged, capable
of navigating codebases, writing code, running tests, and debugging:

\paragraph{Devin and OpenHands.}
Devin (Cognition Labs, 2024) \citep{cognition2024devin} demonstrated end-to-end
autonomous task completion, while OpenHands \citep{wang2024openhands} provides
an open-source alternative with similar capabilities. Both excel at single-agent
execution but focus on individual task autonomy rather than multi-agent
orchestration. SPOQ complements such systems by providing the coordination
layer for parallel execution across multiple agents.

\subsection{Code Generation and Pair Programming Tools}

Developer-facing code generation tools occupy a different point in the
autonomy spectrum, emphasizing human-AI collaboration over full autonomy:

\paragraph{Aider and GPT-Engineer.}
Aider \citep{aider2024} enables conversational code editing where developers
describe changes in natural language and the tool applies them to local
repositories. GPT-Engineer \citep{gpteng2023} generates entire applications
from specifications, producing directory structures and boilerplate code.
Both tools operate as single agents without parallel execution capabilities
or formal validation beyond user inspection.

\paragraph{Claude Code.}
Claude Code (Anthropic, 2025) marked an inflection point for AI-assisted
software development, establishing the agentic coding paradigm where
developers collaborate with an autonomous agent that navigates codebases,
executes commands, and iterates on solutions within a persistent terminal
session. Unlike completion-based tools, Claude Code operates with full
project context and can orchestrate multi-step workflows autonomously,
making it the catalyst that brought agentic development to mainstream
adoption. SPOQ builds on this foundation: where Claude Code enables a
single developer-agent partnership, SPOQ coordinates multiple such agents
in parallel with structured quality gates.

\paragraph{GitHub Copilot.}
GitHub Copilot \citep{chen2021codex} suggests code completions while humans
retain editing control, inverting the SPOQ model: humans perform the work
while AI assists, rather than agents performing work while humans validate.
SPOQ targets a different use case where agents drive execution with human
oversight, enabling higher throughput on parallelizable tasks.

\paragraph{Positioning SPOQ.}
We note that Claude Code, Devin, and similar products are evolving rapidly;
the capabilities described reflect their state as of late 2025.
SPOQ does not compete with execution-focused agents or code generation tools.
Rather, it provides an orchestration layer that coordinates multiple agents
on complex projects. The key insight is that orchestration and execution are
separable concerns: an agent like Claude Code excels at executing individual
tasks, while SPOQ excels at decomposing projects into tasks, scheduling them
for parallel execution, and validating quality at each gate.
